% Cover letter for the submission to Universal Access in the Information Society
\documentclass[11pt,a4paper]{article}
\usepackage[utf8]{inputenc}
\usepackage[T1]{fontenc}
\usepackage{mathptmx}
\usepackage{microtype}
\usepackage[a4paper,top=1.5cm,bottom=1.4cm,left=2.2cm,right=2.2cm]{geometry}
\usepackage{xcolor}
\usepackage[hidelinks]{hyperref}
\definecolor{uteq}{RGB}{31,59,115}
\setlength{\parindent}{0pt}
\setlength{\parskip}{0.45em}
\pagestyle{empty}
\linespread{0.96}

\begin{document}

\begin{flushright}
	\small
	{\color{uteq}\textbf{Gleiston Guerrero-Ulloa, PhD}}\\
	Facultad de Ciencias de la Computaci\'on, Universidad T\'ecnica Estatal de Quevedo\\
	Quevedo, Los R\'ios, Ecuador\\
	\href{mailto:gguerrero@uteq.edu.ec}{gguerrero@uteq.edu.ec} \quad ORCID 0000-0001-5990-2357\\[0.5em]
	23 October 2026
\end{flushright}

\vspace{0.1em}
Prof.\ Margherita Antona and Prof.\ Constantine Stephanidis\\
Editors-in-Chief, \textit{Universal Access in the Information Society}, Springer Nature

\vspace{0.1em}
\textbf{Submission of the manuscript ``Two barriers to equitable access: WCAG conformance and pre-consent tracking on Ecuadorian and top-ranked university websites''}

Dear Professor Antona and Professor Stephanidis,

On behalf of my co-authors, Andrea Z\'u\~niga Paredes and Efra\'in D\'iaz-Mac\'ias, I submit the above manuscript for consideration as an original research article in \textit{Universal Access in the Information Society}.

A university website is the channel through which applicants find programmes, apply and enrol, so any barrier embedded in it restricts access to the institution itself. The manuscript measures two such barriers on the same pages: conformance with the Web Content Accessibility Guidelines (WCAG), and the tracking that begins before a visitor has given any consent. It also asks whose location determines the protection a visitor receives, a question only a multi-location design can answer.

We audited the home page and privacy notice of 126 universities in August 2026, the complete census of the 63 Ecuadorian institutions and a benchmark of 63 institutions drawn from three international rankings, and we replicated the live audit from Ecuador, Germany, the United Kingdom and the United States. Ecuadorian institutions published a privacy notice in 61.9\% of cases, against 93.7\% in the benchmark group, and their accessibility failures concentrated at the most basic tier, with 45\% of failing nodes at level A against 19\%. Pre-consent tracking did not differ between the groups, but it did differ by the visitor's location: on the same day, 67.8\% of sites tracked visitors from Ecuador and 60.3\% from Germany (Holm-adjusted $p=0.023$). The difference was produced by advertising vendors rather than by the universities. Microsoft and TapAd set no cookies for German or British visitors while setting them for Ecuadorian and United States ones, on the same pages.

The work adds three things to the literature the journal has built. It measures accessibility and privacy together, on one set of sites and under one protocol, which the two research traditions have not done. It brings a Latin American vantage point into a geo-differential literature built almost entirely on the axis between Europe and the United States, and it attributes the difference to the named third parties whose cookies change with location, which determines whether a remedy addressed to universities can work at all. And it rests on a complete national population, the one a national regulator must supervise, rather than on a ranked or convenience sample.

The results are also directly usable. For institutions, the most frequent level-A failure, missing accessible names for links, is inexpensive to correct, and publishing a privacy contact closes the largest documentary gap. For policy, the finding that pre-consent protection follows perimeters set by each vendor, rather than the law that binds the institution, matters for every applicant outside those perimeters.

We have tried to make every claim checkable. Each vantage point was measured in three passes and its location was verified 84 times; the exclusions and sensitivity analyses are reported in full; three WCAG success criteria were validated manually with two independent codings; and the limits of the design are stated rather than implied. The raw data, the audit instrument and the analysis scripts are openly available on Zenodo (\href{https://doi.org/10.5281/zenodo.22405998}{10.5281/zenodo.22405998}), and every figure in the article can be regenerated from them.

The manuscript is original, has not been published elsewhere and is not under consideration by any other journal. All authors have read and approved it. The study received no funding. Two authors are affiliated with one of the audited institutions; this is declared under competing interests, with the safeguards applied. The use of a generative artificial intelligence assistant is declared in the manuscript, including the one point at which it reached the data.

Thank you for considering our work. We look forward to the reviewers' comments.

\noindent\begin{minipage}{\linewidth}
	Yours sincerely,\\[1.4em]
	Gleiston Guerrero-Ulloa, PhD\\
	Corresponding author, on behalf of all authors
\end{minipage}

\end{document}
